\documentclass{symbulator}
\booktitle{How to use this}
\booksubtitle{Editing, building and publishing the Symbulator documentation}
\bookversion{The documentation toolkit}
\bookplatform{One source tree, three PDFs, one website}
\bookauthor{Written for Roberto Perez-Franco}

\begin{document}
\maketitlepage
\tableofcontents

\frontchapter{Where everything lives}
\chaptersummary{Unpack the zip once, and that folder becomes the source of
truth for every version of the documentation. Everything else is generated
from it.}

Download \symcode{symbulator-docs.zip} and unpack it wherever you keep your
projects. Inside you will find this:

\begin{itemize}
\item \symcode{src/} --- the chapters, one Markdown file each. \textbf{This is
      what you edit.}
\item \symcode{book.yaml} --- chapter order, the three versions, and the words
      that change between them.
\item \symcode{tex/symbulator.cls} --- the look of the PDFs.
\item \symcode{web/} --- the look of the website: \symcode{index.php},
      \symcode{assets/style.css}, \symcode{.htaccess}.
\item \symcode{assets/} --- the figures.
\item \symcode{tools/} --- helper scripts.
\item \symcode{build/} --- generated output. \textbf{Never edit anything in
      here}; the next build overwrites it.
\end{itemize}

\begin{calloutwarning}{The copy in the zip is a snapshot}
The \symcode{build/} folder inside the zip is there so you can look at the
results straight away. The moment you run a build it is replaced.
\end{calloutwarning}

\frontchapter{What you need installed}
\chaptersummary{Python for everything, and XeLaTeX only if you want the PDFs.
PHP is not needed on your machine at all.}

\section{Python}

Python 3.9 or newer. From inside the project folder, install the two libraries
the build uses:

\begin{symtype}
pip install -r requirements.txt
\end{symtype}

That is PyYAML, which reads \symcode{book.yaml} and the front matter of each
chapter, and Matplotlib, which is only used by the figure-placeholder script.

\section{XeLaTeX, for the PDFs}

The PDFs are typeset with XeLaTeX. On Windows, install MiKTeX and let it fetch
packages on demand the first time you build. On macOS, install MacTeX. On
Linux, install \symcode{texlive-xetex}, \symcode{texlive-fonts-recommended} and
\symcode{texlive-latex-extra}.

If you skip this, everything still works except the PDFs: build with
\symcode{-{}-web} and you get the website only.

\section{PHP}

You do not need PHP on your own machine. It runs on the server, and your cPanel
hosting already has it. To preview the site locally you use the static preview
described later, which needs nothing but a browser.

\frontchapter{Editing the source}
\chaptersummary{The chapters are plain Markdown with a small set of additions.
Any text editor will open them, but a code editor plus Claude Code will save
you a great deal of typing.}

\section{What to use}

I would use \textbf{VS Code}, with \textbf{Claude Code} running in the same
folder. You already use that pairing for the Android app, and it pays off more
here: because every chapter is plain UTF-8 text in one tree, you can ask for
sweeping changes --- ``convert Lesson 10 from this HTML page'', ``fix the
version 9 code everywhere it appears'' --- and have them applied across all
fifteen files at once.

VS Code also gives you syntax colouring inside the fenced code blocks, and a
project-wide search, which is how you will work through the outstanding items:

\begin{symtype}
grep -rn TODO src/
\end{symtype}

\begin{calloutnote}{The built-in Markdown preview will disappoint you}
VS Code's preview pane does not understand the \symcode{:::} directives or the
\symcode{\{\{v7|...\}\}} spans, and shows them as literal text. The real preview
is the generated one, described in the next chapter.
\end{calloutnote}

\section{The markup, in one page}

Ordinary Markdown handles headings, lists, bold, italic and links. On top of
that:

\begin{itemize}
\item \symcode{\#\# Title} is a numbered section; \symcode{\#\#\# Title} is a
      subsection.
\item \symcode{\{\{v7|text\}\}} shows text only in version 7;
      \symcode{\{\{v7,8|text\}\}} in two of them; \symcode{\{\{!v8|text\}\}} in
      everything except version 8.
\item \symcode{\{\{t:container\}\}} inserts a word that changes with the version
      --- folder, document or session --- defined in \symcode{book.yaml}.
\item \symcode{\{\{i:resistor\}\}} adds an index entry, invisibly.
\item \symcode{\{\{ref:lesson-dc\}\}} cross-references another chapter or
      section, and renders as ``Lesson 1'' or ``section 1.3'' with the right
      number for that version.
\item \symcode{:::~tip Title} opens a callout; also \symcode{note},
      \symcode{warning} and \symcode{danger}. Close every one with
      \symcode{:::} on its own line.
\item \symcode{:::~problem}, \symcode{:::~answer}, \symcode{:::~figure} and
      \symcode{:::~practice} build the worked examples.
\item \symcode{:::~only 7,8} and \symcode{:::~not 8} gate whole blocks.
\item A fence marked \symcode{sym 7} is something the reader types on version 7;
      \symcode{out} is something the software returns. Consecutive fences with
      different version tags are variants of the same instruction.
\end{itemize}

The full specification is in \symcode{SPEC.md} in the project folder.

\section{The standfirst}

Each chapter opens with a larger paragraph telling the reader what to expect.
It comes from \symcode{summary:} in that chapter's front matter, and it is the
same text that appears on the landing page card.

\frontchapter{Building}
\chaptersummary{One command produces all three PDFs and the whole website. A
second one watches the folder and rebuilds as you type.}

\section{Check before you build}

\begin{symtype}
python3 build.py -{}-check
\end{symtype}

This validates every cross-reference, every version term, every code-fence
version tag and every figure, for all three versions, and tells you what is
wrong. It also lists the figures still to be supplied. The same check runs at
the start of every build, as a warning.

\section{Build everything}

\begin{symtype}
python3 build.py
\end{symtype}

That writes:

\begin{itemize}
\item \symcode{build/pdf/} --- \symcode{symbulator-v7.pdf},
      \symcode{-v8.pdf} and \symcode{-v9.pdf}
\item \symcode{build/web/} --- the whole site, ready to upload
\item \symcode{build/tex/} --- the generated LaTeX, kept for when something
      goes wrong
\end{itemize}

Useful variations: \symcode{-{}-web} for the website only, which is fast;
\symcode{-{}-pdf} for the PDFs only; \symcode{-{}-versions 8} to work on one
version at a time.

\section{While you write}

\begin{symtype}
python3 tools/watch.py
\end{symtype}

Leave this running in a terminal. It rebuilds about a second after you save a
file in \symcode{src/}, and refreshes the preview. Open
\symcode{build/preview/v7-home.html} in a browser and reload the tab as you go.
Stop it with Ctrl-C.

It skips the PDFs by default, because XeLaTeX takes roughly a minute for all
three. Add \symcode{-{}-pdf} when you want those rebuilt too --- usually only
before you publish.

\begin{callouttip}{The preview is not the site}
\symcode{build/preview/} is flat HTML, generated so you can look at the design
without a PHP server. What you upload is \symcode{build/web/}.
\end{callouttip}

\frontchapter{Figures}
\chaptersummary{Every schematic needs two files with the same name: an SVG for
the website and a PDF for print. Until the real artwork exists, the build draws
a labelled placeholder.}

Put figures in \symcode{assets/circuit/}, as \symcode{name.svg} and
\symcode{name.pdf}. The source refers to the SVG; the LaTeX side drops the
extension so XeLaTeX picks up the PDF automatically.

For anything referred to but not yet drawn:

\begin{symtype}
python3 tools/make\_placeholders.py
\end{symtype}

This draws a dashed box carrying the caption and the filename, in both formats,
so the build stays green and every gap announces itself in the document. Drop
the real artwork over a placeholder and the script will not touch it again.

\frontchapter{Publishing}
\chaptersummary{One upload serves all three versions. There are no build tools
on the server.}

Copy the contents of \symcode{build/web/} into the documentation folder on your
cPanel hosting, and put the three PDFs beside \symcode{index.php} so the
download links resolve.

That is the whole deployment. \symcode{index.php} reads
\symcode{content/v7}, \symcode{v8} or \symcode{v9} according to the
\symcode{?v=} in the URL, so the version selector works server-side with no
JavaScript. The bundled \symcode{.htaccess} turns that into tidier addresses
such as \symcode{/7/lesson-dc}.

\begin{calloutnote}{If the pretty URLs do not work}
Some cPanel setups have \symcode{AllowOverride} switched off, which makes
Apache ignore \symcode{.htaccess}. The plain \symcode{?v=7\&p=lesson-dc} form
keeps working regardless, so nothing breaks --- the addresses are just longer.
\end{calloutnote}

\frontchapter{When something goes wrong}
\chaptersummary{Three failures account for almost everything, and each one
announces itself clearly.}

\section{An unclosed directive}

If the build stops with \emph{unclosed \symcode{:::} directive}, a block was
opened and never closed. Every \symcode{:::~tip}, \symcode{:::~problem} and
\symcode{:::~only} needs its own \symcode{:::} on a line by itself, including
the nested ones.

\section{A broken cross-reference}

\symcode{-{}-check} reports the chapter, the label and the version. Either the
label is misspelled, or it points at a chapter that does not exist in that
version --- worth thinking about, because it may mean the sentence should be
version-tagged too.

\section{XeLaTeX stops}

Read \symcode{build/tex/symbulator-v7.log} and search for a line starting with
an exclamation mark. The generated \symcode{.tex} is kept beside it so you can
see exactly what was fed to LaTeX.

\end{document}
